% ID: SEL-D0030-OA-E02
% ESTADO: validada. S031, cotejada con examen y criterios.
\textbf{a)} $f'(x)=2x-6$ es negativa para $0<x<3$ y
positiva para $x>3$. El coste disminuye en $[0,3]$ y luego
aumenta. Su mínimo es $\boxed{31}$ al producir $\boxed{3}$ unidades.

\textbf{b)} Sin producción, $f(0)=\boxed{40}$.
La ecuación $f(x)=80$ da $x^2-6x-40=(x-10)(x+4)=0$.
Por el dominio $x\ge0$, la producción es de $\boxed{10}$ unidades.

\textbf{c)} Como $f(x)=(x-3)^2+31$, la gráfica es la parte
de una parábola convexa para $x\ge0$. Pasa por $(0,40)$,
tiene vértice $(3,31)$ y pasa por $(10,80)$.
\begin{center}
\setlength{\unitlength}{6mm}
\begin{picture}(15,10)(-2,-1)
\qbezier(0,0)(5.5,0)(11,0)\put(11.2,0){$x$}
\qbezier(0,0)(0,4.3)(0,8.6)\put(0,8.7){$f$}
\qbezier(0,4)(5,1)(10,8)
\put(-0.2,4){\makebox(0,0)[r]{$(0,40)$}}
\put(3,2.9){\makebox(0,0)[t]{$(3,31)$}}
\put(10.2,8){$(10,80)$}
\end{picture}
\end{center}
